\section{Protokol Diffie-Hellman}

Protokol Diffie-Hellman (1976) memungkinkan dua pihak menetapkan kunci rahasia bersama melalui saluran tidak aman \cite{diffie1976,katz2020,trappe2006}. Berdasarkan masalah diskret logaritma.

\begin{figure}[!htbp]
\centering
\begin{tikzpicture}[
  node distance=1.2cm,
  box/.style={rectangle, draw, minimum width=2cm, minimum height=0.6cm, align=center, font=\small}
]
\node[box, fill=blue!15] (alice) {Alice};
\node[box, fill=green!15, right=3.5cm of alice] (bob) {Bob};
\draw[-Stealth] (alice) -- node[above, font=\scriptsize] {$A = g^a \bmod p$} (bob);
\draw[-Stealth] (bob) -- node[below, font=\scriptsize] {$B = g^b \bmod p$} (alice);
\node[below=0.8cm of $(alice)!0.5!(bob)$, font=\scriptsize] {Kunci bersama: $K = g^{ab} \bmod p$};
\end{tikzpicture}
\caption{Protokol Diffie-Hellman: pertukaran nilai untuk menghasilkan kunci rahasia}
\label{fig:dh}
\end{figure}

Alice dan Bob menyepakati bilangan prima $p$ dan generator $g$. Alice memilih $a$, mengirim $A = g^a \bmod p$. Bob memilih $b$, mengirim $B = g^b \bmod p$. Keduanya menghitung $K = g^{ab} \bmod p$ (Alice: $B^a$, Bob: $A^b$). Eve yang menyadap $A$ dan $B$ tidak dapat menghitung $K$ tanpa menyelesaikan diskret logaritma \cite{diffie1976,katz2020,trappe2006}.

\begin{lstlisting}[style=matlab,caption={Diffie-Hellman dalam MATLAB}]
function K = diffie_hellman(p, g, a, b)
    A = mod_pow(g, a, p);  % Alice mengirim A
    B = mod_pow(g, b, p);  % Bob mengirim B
    K_alice = mod_pow(B, a, p);  % K = B^a mod p
    K_bob = mod_pow(A, b, p);    % K = A^b mod p
    K = K_alice;  % K_alice == K_bob
end
\end{lstlisting}
